
\documentclass{article}
\usepackage{colortbl}
\usepackage{makecell}
\usepackage{multirow}
\usepackage{supertabular}

\begin{document}

\newcounter{utterance}

\centering \large Interaction Transcript for game `clean\_up', experiment `0\_easy\_5obj\_tr', episode 3 with nemotron{-}3{-}ultra.
\vspace{24pt}

{ \footnotesize  \setcounter{utterance}{1}
\setlength{\tabcolsep}{0pt}
\begin{supertabular}{c@{$\;$}|p{.15\linewidth}@{}p{.15\linewidth}p{.15\linewidth}p{.15\linewidth}p{.15\linewidth}p{.15\linewidth}}
    \# & $\;$A & \multicolumn{4}{c}{Game Master} & $\;\:$B\\
    \hline

    \theutterance \stepcounter{utterance}  
    & & \multicolumn{4}{p{0.6\linewidth}}{
        \cellcolor[rgb]{0.9,0.9,0.9}{
            \makecell[{{p{\linewidth}}}]{
                \texttt{\tiny{[P1$\langle$GM]}}
                \texttt{Ben senin oyun yöneticinim ve aşağıdaki ızgarayı oyun tahtası olarak kullandığınız işbirlikçi bir oyun oynuyorsunuz.} \\
\texttt{```} \\
\texttt{ 1234567} \\
\texttt{╔════╤══╗} \\
\texttt{║H◌◌◌│◌◌║ 1} \\
\texttt{║◌W◌◌│◌◌║ 2} \\
\texttt{║◌C◌◌│◌◌║ 3} \\
\texttt{║T◌◌◌│I◌║ 4} \\
\texttt{╟──┐◌│◌◌║ 5} \\
\texttt{║◌◌├─┼──╢ 6} \\
\texttt{║◌◌│◌│◌◌║ 7} \\
\texttt{╚══╧═╧══╝} \\
\\ 
\texttt{```} \\
\\ 
\texttt{* Üst kenarda x koordinatları sağa doğru artar, sağ kenarda ise y koordinatları aşağı doğru artar.} \\
\texttt{* Aşağıdaki nesneler oyun tahtasına rastgele yerleştirilmiştir: 'W', 'I', 'T', 'C', 'H'.} \\
\\ 
\\ 
\texttt{Diğer oyuncu, nesnelerin farklı rastgele konumlara yerleştirildiği oyun tahtasının bir varyasyonunu görür. Sen diğer oyuncunun tahtasını göremezsin ve onlar da seninkini göremez.} \\
\texttt{**Amaç:**} \\
\\ 
\texttt{Her iki oyuncunun da, kendi oyun tahtalarındaki nesneleri hareket ettirerek aynı türden nesnelerin aynı koordinatlara gelmesini sağlaması gerekir. Ortak bir hedef düzen üzerinde anlaşmak için diğer oyuncuyla iletişim kurmalısınız.} \\
\\ 
\\ 
\texttt{**Kurallar:**} \\
\\ 
\texttt{* Her turda, aşağıdaki iki komuttan yalnızca birini gönderebilirsin:} \\
\texttt{ 1. `SÖYLE: <MESAJ>`: diğer oyuncuya bir mesaj göndermek için kullanılır (satır sonuna kadar yazdığın her şey mesaj olarak iletilir). Ben bu mesajı partnerine ileteceğim.} \\
\texttt{2. `HAMLE: <NESNE>, (<X>, <Y>)`: bir nesneyi yeni bir konuma taşımak için kullanılır. Burada <X> sütunu, <Y> ise satırı ifade eder. Yaptığın hamlenin başarılı olup olmadığını sana söyleyeceğim.} \\
\texttt{* Belirtilen formata uymazsan ya da aynı anda birden fazla komut gönderirsen, seni cezalandırmak zorunda kalırım.} \\
\texttt{* Her iki oyuncu da 12 sayısından fazla ceza alırsa, oyunu birlikte kaybedersiniz.} \\
\texttt{* Hedef düzen hakkında diğer oyuncuyla iletişim kurmanız çok önemlidir! Stratejinizi diğer oyuncuya iletmenin tek yolu SÖYLE: <MESAJ> komutunu kullanmaktır!} \\
\\ 
\\ 
\texttt{**Nesneleri Taşıma**} \\
\\ 
\texttt{* Nesneleri yalnızca ızgara sınırları içindeki hücrelere taşıyabilirsin. Hedef hücre boş olmalıdır; yani yalnızca ‘◌’ sembolünü içermelidir.} \\
\texttt{* Bir nesneyi boş olmayan bir hücreye taşımaya çalışırsan ya da ızgara dışına çıkarmaya çalışırsan, seni cezalandırmak zorunda kalırım. Ardından tekrar deneme hakkın olur.} \\
\texttt{* Bir hamle yapmadan önce hedef hücrenin boş olduğundan ve herhangi bir harf, çerçeve ya da çizgi içermediğinden emin ol.} \\
\\ 
\\ 
\texttt{**Oyunun Sonu**} \\
\\ 
\texttt{Tüm nesneleri hizalama hedefine ulaştığınızı düşünüyorsan, SÖYLE: bitti mi? göndererek diğer oyuncudan oyunu bitirmesini isteyebilirsin. Eğer diğer oyuncu senden oyunu bitirmeni isterse ve sen SÖYLE: bitti! ile yanıt verirsen, oyun sona erer.} \\
\\ 
\texttt{Oyun, 20 tur içinde sona ererse her iki oyuncu da kazanır. Bir tur, iki oyuncunun da geçerli bir komut göndermesi olarak tanımlanır.} \\
\\ 
\texttt{**Puanlama:**} \\
\\ 
\texttt{Aynı nesneler her iki oyun tahtasında ne kadar yakın konumlanırsa, o kadar fazla puan kazanırsınız. Cezalar puanınızı düşürür. Rekoru kırabilir misiniz?} \\
\\ 
\\ 
\texttt{Oyuna başlamak için diğer oyuncuya bir mesaj gönder!} \\
            }
        }
    }
    & \\ \\

    \theutterance \stepcounter{utterance}  
    & \multicolumn{4}{p{0.6\linewidth}}{
        \cellcolor[rgb]{0.9,0.9,0.9}{
            \makecell[{{p{\linewidth}}}]{
                \texttt{\tiny{[P1$\rangle$GM]}}
                \texttt{SÖYLE: Merhaba! Oyuna hoş geldin. Benim tahtamda nesneler şu koordinatlarda: H(1,1), W(2,2), C(2,3), T(1,4), I(5,4). Sende nerelerde? Hedef düzen için bir fikrin var mı?} \\
            }
        }
    }
    & & \\ \\

    \theutterance \stepcounter{utterance}  
    & & & \multicolumn{2}{p{0.3\linewidth}}{
        \cellcolor[rgb]{0.9,0.9,0.9}{
            \makecell[{{p{\linewidth}}}]{
                \texttt{\tiny{[GM$|$GM]}}
                \texttt{SÖYLE: Merhaba! Oyuna hoş geldin. Benim tahtamda nesneler şu koordinatlarda: H(1,1), W(2,2), C(2,3), T(1,4), I(5,4). Sende nerelerde? Hedef düzen için bir fikrin var mı?} \\
            }
        }
    }
    & & \\ \\

    \theutterance \stepcounter{utterance}  
    & & \multicolumn{4}{p{0.6\linewidth}}{
        \cellcolor[rgb]{0.9,0.9,0.9}{
            \makecell[{{p{\linewidth}}}]{
                \texttt{\tiny{[GM$\rangle$P2]}}
                \texttt{Ben senin oyun yöneticinim ve aşağıdaki ızgarayı oyun tahtası olarak kullandığınız işbirlikçi bir oyun oynuyorsunuz.} \\
\texttt{```} \\
\texttt{ 1234567} \\
\texttt{╔════╤══╗} \\
\texttt{║◌◌◌T│◌◌║ 1} \\
\texttt{║◌◌◌◌│◌◌║ 2} \\
\texttt{║◌◌◌◌│◌◌║ 3} \\
\texttt{║W◌◌◌│◌◌║ 4} \\
\texttt{╟──┐◌│◌◌║ 5} \\
\texttt{║H◌├─┼──╢ 6} \\
\texttt{║I◌│◌│◌C║ 7} \\
\texttt{╚══╧═╧══╝} \\
\\ 
\texttt{```} \\
\\ 
\texttt{* Üst kenarda x koordinatları sağa doğru artar, sağ kenarda ise y koordinatları aşağı doğru artar.} \\
\texttt{* Aşağıdaki nesneler oyun tahtasına rastgele yerleştirilmiştir: 'W', 'I', 'T', 'C', 'H'.} \\
\\ 
\\ 
\texttt{Diğer oyuncu, nesnelerin farklı rastgele konumlara yerleştirildiği oyun tahtasının bir varyasyonunu görür. Sen diğer oyuncunun tahtasını göremezsin ve onlar da seninkini göremez.} \\
\texttt{**Amaç:**} \\
\\ 
\texttt{Her iki oyuncunun da, kendi oyun tahtalarındaki nesneleri hareket ettirerek aynı türden nesnelerin aynı koordinatlara gelmesini sağlaması gerekir. Ortak bir hedef düzen üzerinde anlaşmak için diğer oyuncuyla iletişim kurmalısınız.} \\
\\ 
\\ 
\texttt{**Kurallar:**} \\
\\ 
\texttt{* Her turda, aşağıdaki iki komuttan yalnızca birini gönderebilirsin:} \\
\texttt{ 1. `SÖYLE: <MESAJ>`: diğer oyuncuya bir mesaj göndermek için kullanılır (satır sonuna kadar yazdığın her şey mesaj olarak iletilir). Ben bu mesajı partnerine ileteceğim.} \\
\texttt{2. `HAMLE: <NESNE>, (<X>, <Y>)`: bir nesneyi yeni bir konuma taşımak için kullanılır. Burada <X> sütunu, <Y> ise satırı ifade eder. Yaptığın hamlenin başarılı olup olmadığını sana söyleyeceğim.} \\
\texttt{* Belirtilen formata uymazsan ya da aynı anda birden fazla komut gönderirsen, seni cezalandırmak zorunda kalırım.} \\
\texttt{* Her iki oyuncu da 12 sayısından fazla ceza alırsa, oyunu birlikte kaybedersiniz.} \\
\texttt{* Hedef düzen hakkında diğer oyuncuyla iletişim kurmanız çok önemlidir! Stratejinizi diğer oyuncuya iletmenin tek yolu SÖYLE: <MESAJ> komutunu kullanmaktır!} \\
\\ 
\\ 
\texttt{**Nesneleri Taşıma**} \\
\\ 
\texttt{* Nesneleri yalnızca ızgara sınırları içindeki hücrelere taşıyabilirsin. Hedef hücre boş olmalıdır; yani yalnızca ‘◌’ sembolünü içermelidir.} \\
\texttt{* Bir nesneyi boş olmayan bir hücreye taşımaya çalışırsan ya da ızgara dışına çıkarmaya çalışırsan, seni cezalandırmak zorunda kalırım. Ardından tekrar deneme hakkın olur.} \\
\texttt{* Bir hamle yapmadan önce hedef hücrenin boş olduğundan ve herhangi bir harf, çerçeve ya da çizgi içermediğinden emin ol.} \\
\\ 
\\ 
\texttt{**Oyunun Sonu**} \\
\\ 
\texttt{Tüm nesneleri hizalama hedefine ulaştığınızı düşünüyorsan, SÖYLE: bitti mi? göndererek diğer oyuncudan oyunu bitirmesini isteyebilirsin. Eğer diğer oyuncu senden oyunu bitirmeni isterse ve sen SÖYLE: bitti! ile yanıt verirsen, oyun sona erer.} \\
\\ 
\texttt{Oyun, 20 tur içinde sona ererse her iki oyuncu da kazanır. Bir tur, iki oyuncunun da geçerli bir komut göndermesi olarak tanımlanır.} \\
\\ 
\texttt{**Puanlama:**} \\
\\ 
\texttt{Aynı nesneler her iki oyun tahtasında ne kadar yakın konumlanırsa, o kadar fazla puan kazanırsınız. Cezalar puanınızı düşürür. Rekoru kırabilir misiniz?} \\
\\ 
\texttt{Diğer oyuncu şu mesajı göndererek oyunu başlattı:} \\
\\ 
\texttt{"Merhaba! Oyuna hoş geldin. Benim tahtamda nesneler şu koordinatlarda: H(1,1), W(2,2), C(2,3), T(1,4), I(5,4). Sende nerelerde? Hedef düzen için bir fikrin var mı?"} \\
\\ 
\texttt{İlk komutun ne?} \\
            }
        }
    }
    & \\ \\

    \theutterance \stepcounter{utterance}  
    & & & \multicolumn{4}{p{0.6\linewidth}}{
        \cellcolor[rgb]{0.9,0.9,0.9}{
            \makecell[{{p{\linewidth}}}]{
                \texttt{\tiny{[GM$\langle$P2]}}
                \texttt{SÖYLE: Merhaba! Benim tahtamda nesneler şu koordinatlarda: T(4,1), W(1,4), H(1,6), I(1,7), C(7,7). Hedef düzen için köşeleri kullanmak mantıklı görünüyor. H ve W sol altta (H(1,6), W(1,5)), I ve C sağ altta (I(7,6), C(7,7)), T ise ortada (T(4,1)) gibi bir yerleşim öneririm. Bu düzen sana uygun mu?} \\
            }
        }
    }
     \\ \\

    \theutterance \stepcounter{utterance}  
    & & & \multicolumn{2}{p{0.3\linewidth}}{
        \cellcolor[rgb]{0.9,0.9,0.9}{
            \makecell[{{p{\linewidth}}}]{
                \texttt{\tiny{[GM$|$GM]}}
                \texttt{SÖYLE: Merhaba! Benim tahtamda nesneler şu koordinatlarda: T(4,1), W(1,4), H(1,6), I(1,7), C(7,7). Hedef düzen için köşeleri kullanmak mantıklı görünüyor. H ve W sol altta (H(1,6), W(1,5)), I ve C sağ altta (I(7,6), C(7,7)), T ise ortada (T(4,1)) gibi bir yerleşim öneririm. Bu düzen sana uygun mu?} \\
            }
        }
    }
    & & \\ \\

    \theutterance \stepcounter{utterance}  
    & & \multicolumn{4}{p{0.6\linewidth}}{
        \cellcolor[rgb]{0.9,0.9,0.9}{
            \makecell[{{p{\linewidth}}}]{
                \texttt{\tiny{[P1$\langle$GM]}}
                \texttt{Mesajın diğer oyuncuya iletildi.} \\
\\ 
\\ 
\texttt{Şu anda en fazla 20 turdan 1. turu oynuyorsunuz.} \\
\texttt{Toplamda 12 cezanın 0 kadarını biriktirdiniz.Diğer oyuncu şu mesajı gönderdi:} \\
\\ 
\\ 
\texttt{"Merhaba! Benim tahtamda nesneler şu koordinatlarda: T(4,1), W(1,4), H(1,6), I(1,7), C(7,7). Hedef düzen için köşeleri kullanmak mantıklı görünüyor. H ve W sol altta (H(1,6), W(1,5)), I ve C sağ altta (I(7,6), C(7,7)), T ise ortada (T(4,1)) gibi bir yerleşim öneririm. Bu düzen sana uygun mu?"} \\
\\ 
\\ 
\texttt{Bir sonraki komutun nedir?} \\
            }
        }
    }
    & \\ \\

    \theutterance \stepcounter{utterance}  
    & \multicolumn{4}{p{0.6\linewidth}}{
        \cellcolor[rgb]{0.9,0.9,0.9}{
            \makecell[{{p{\linewidth}}}]{
                \texttt{\tiny{[P1$\rangle$GM]}}
                \texttt{SÖYLE: Öneri iyi, köşeler sabitlenmesi kolay. Benim tahtamda 5. satır ve 6. satırda duvarlar (┐, ├, ─, ┼, │, ┘) var, bu hücreler engel. H(1,6) ve W(1,5) hedefleri boş ve erişilebilir görünüyor. I(7,6), C(7,7) ve T(4,1) de boş. Bu hedef düzen kabul ediyorum. Ben T'yi (4,1)'e taşımaya başlıyorum. Sen de kendi T'ni (4,1)'e taşı.} \\
            }
        }
    }
    & & \\ \\

    \theutterance \stepcounter{utterance}  
    & & & \multicolumn{2}{p{0.3\linewidth}}{
        \cellcolor[rgb]{0.9,0.9,0.9}{
            \makecell[{{p{\linewidth}}}]{
                \texttt{\tiny{[GM$|$GM]}}
                \texttt{SÖYLE: Öneri iyi, köşeler sabitlenmesi kolay. Benim tahtamda 5. satır ve 6. satırda duvarlar (┐, ├, ─, ┼, │, ┘) var, bu hücreler engel. H(1,6) ve W(1,5) hedefleri boş ve erişilebilir görünüyor. I(7,6), C(7,7) ve T(4,1) de boş. Bu hedef düzen kabul ediyorum. Ben T'yi (4,1)'e taşımaya başlıyorum. Sen de kendi T'ni (4,1)'e taşı.} \\
            }
        }
    }
    & & \\ \\

    \theutterance \stepcounter{utterance}  
    & & \multicolumn{4}{p{0.6\linewidth}}{
        \cellcolor[rgb]{0.9,0.9,0.9}{
            \makecell[{{p{\linewidth}}}]{
                \texttt{\tiny{[GM$\rangle$P2]}}
                \texttt{Mesajın diğer oyuncuya iletildi.} \\
\\ 
\\ 
\texttt{Şu anda en fazla 20 turdan 2. turu oynuyorsunuz.} \\
\texttt{Toplamda 12 cezanın 0 kadarını biriktirdiniz.Diğer oyuncu şu mesajı gönderdi:} \\
\\ 
\\ 
\texttt{"Öneri iyi, köşeler sabitlenmesi kolay. Benim tahtamda 5. satır ve 6. satırda duvarlar (┐, ├, ─, ┼, │, ┘) var, bu hücreler engel. H(1,6) ve W(1,5) hedefleri boş ve erişilebilir görünüyor. I(7,6), C(7,7) ve T(4,1) de boş. Bu hedef düzen kabul ediyorum. Ben T'yi (4,1)'e taşımaya başlıyorum. Sen de kendi T'ni (4,1)'e taşı."} \\
\\ 
\\ 
\texttt{Bir sonraki komutun nedir?} \\
            }
        }
    }
    & \\ \\

    \theutterance \stepcounter{utterance}  
    & & & \multicolumn{4}{p{0.6\linewidth}}{
        \cellcolor[rgb]{0.9,0.9,0.9}{
            \makecell[{{p{\linewidth}}}]{
                \texttt{\tiny{[GM$\langle$P2]}}
                \texttt{HAMLE: T, (4, 1)} \\
            }
        }
    }
     \\ \\

    \theutterance \stepcounter{utterance}  
    & & & \multicolumn{2}{p{0.3\linewidth}}{
        \cellcolor[rgb]{0.9,0.9,0.9}{
            \makecell[{{p{\linewidth}}}]{
                \texttt{\tiny{[GM$|$GM]}}
                \texttt{HAMLE: T, (4, 1)} \\
            }
        }
    }
    & & \\ \\

    \theutterance \stepcounter{utterance}  
    & & & \multicolumn{2}{p{0.3\linewidth}}{
        \cellcolor[rgb]{0.9,0.9,0.9}{
            \makecell[{{p{\linewidth}}}]{
                \texttt{\tiny{[GM$|$GM]}}
                \texttt{Ceza: (4,1) boş değil; "T" içeriyor.} \\
\\ 
\texttt{Toplamda 12 cezanın 1 kadarını biriktirdiniz.Lütfen tekrar deneyin!} \\
            }
        }
    }
    & & \\ \\

    \theutterance \stepcounter{utterance}  
    & & \multicolumn{4}{p{0.6\linewidth}}{
        \cellcolor[rgb]{0.9,0.9,0.9}{
            \makecell[{{p{\linewidth}}}]{
                \texttt{\tiny{[GM$\rangle$P2]}}
                \texttt{Ceza: (4,1) boş değil; "T" içeriyor.} \\
\\ 
\texttt{Toplamda 12 cezanın 1 kadarını biriktirdiniz.Lütfen tekrar deneyin!} \\
            }
        }
    }
    & \\ \\

    \theutterance \stepcounter{utterance}  
    & & & \multicolumn{4}{p{0.6\linewidth}}{
        \cellcolor[rgb]{0.9,0.9,0.9}{
            \makecell[{{p{\linewidth}}}]{
                \texttt{\tiny{[GM$\langle$P2]}}
                \texttt{SÖYLE: Anlaşıldı, T zaten (4,1) konumundaydı, hamle yapmaya gerek kalmadı. Diğer nesneleri hedef konumlarına taşımaya başlıyorum. Önce W'yi (1,5) konumuna taşıyorum.} \\
            }
        }
    }
     \\ \\

    \theutterance \stepcounter{utterance}  
    & & & \multicolumn{2}{p{0.3\linewidth}}{
        \cellcolor[rgb]{0.9,0.9,0.9}{
            \makecell[{{p{\linewidth}}}]{
                \texttt{\tiny{[GM$|$GM]}}
                \texttt{SÖYLE: Anlaşıldı, T zaten (4,1) konumundaydı, hamle yapmaya gerek kalmadı. Diğer nesneleri hedef konumlarına taşımaya başlıyorum. Önce W'yi (1,5) konumuna taşıyorum.} \\
            }
        }
    }
    & & \\ \\

    \theutterance \stepcounter{utterance}  
    & & \multicolumn{4}{p{0.6\linewidth}}{
        \cellcolor[rgb]{0.9,0.9,0.9}{
            \makecell[{{p{\linewidth}}}]{
                \texttt{\tiny{[P1$\langle$GM]}}
                \texttt{Mesajın diğer oyuncuya iletildi.} \\
\\ 
\\ 
\texttt{Şu anda en fazla 20 turdan 2. turu oynuyorsunuz.} \\
\texttt{Toplamda 12 cezanın 1 kadarını biriktirdiniz.Diğer oyuncu şu mesajı gönderdi:} \\
\\ 
\\ 
\texttt{"Anlaşıldı, T zaten (4,1) konumundaydı, hamle yapmaya gerek kalmadı. Diğer nesneleri hedef konumlarına taşımaya başlıyorum. Önce W'yi (1,5) konumuna taşıyorum."} \\
\\ 
\\ 
\texttt{Bir sonraki komutun nedir?} \\
            }
        }
    }
    & \\ \\

    \theutterance \stepcounter{utterance}  
    & \multicolumn{4}{p{0.6\linewidth}}{
        \cellcolor[rgb]{0.9,0.9,0.9}{
            \makecell[{{p{\linewidth}}}]{
                \texttt{\tiny{[P1$\rangle$GM]}}
                \texttt{HAMLE: T, (4, 1)} \\
            }
        }
    }
    & & \\ \\

    \theutterance \stepcounter{utterance}  
    & & & \multicolumn{2}{p{0.3\linewidth}}{
        \cellcolor[rgb]{0.9,0.9,0.9}{
            \makecell[{{p{\linewidth}}}]{
                \texttt{\tiny{[GM$|$GM]}}
                \texttt{HAMLE: T, (4, 1)} \\
            }
        }
    }
    & & \\ \\

    \theutterance \stepcounter{utterance}  
    & & & \multicolumn{2}{p{0.3\linewidth}}{
        \cellcolor[rgb]{0.9,0.9,0.9}{
            \makecell[{{p{\linewidth}}}]{
                \texttt{\tiny{[GM$|$GM]}}
                \texttt{T moved to (4, 1)} \\
            }
        }
    }
    & & \\ \\

    \theutterance \stepcounter{utterance}  
    & & \multicolumn{4}{p{0.6\linewidth}}{
        \cellcolor[rgb]{0.9,0.9,0.9}{
            \makecell[{{p{\linewidth}}}]{
                \texttt{\tiny{[GM$\rangle$P2]}}
                \texttt{Mesajın diğer oyuncuya iletildi.} \\
\\ 
\\ 
\texttt{Şu anda en fazla 20 turdan 3. turu oynuyorsunuz.} \\
\texttt{Toplamda 12 cezanın 1 kadarını biriktirdiniz.Diğer oyuncu kendi ızgarasında bir nesneyi hareket ettirdi.} \\
\\ 
\\ 
\texttt{Bir sonraki komutun nedir?} \\
            }
        }
    }
    & \\ \\

    \theutterance \stepcounter{utterance}  
    & & & \multicolumn{4}{p{0.6\linewidth}}{
        \cellcolor[rgb]{0.9,0.9,0.9}{
            \makecell[{{p{\linewidth}}}]{
                \texttt{\tiny{[GM$\langle$P2]}}
                \texttt{HAMLE: W, (1, 5)} \\
            }
        }
    }
     \\ \\

    \theutterance \stepcounter{utterance}  
    & & & \multicolumn{2}{p{0.3\linewidth}}{
        \cellcolor[rgb]{0.9,0.9,0.9}{
            \makecell[{{p{\linewidth}}}]{
                \texttt{\tiny{[GM$|$GM]}}
                \texttt{HAMLE: W, (1, 5)} \\
            }
        }
    }
    & & \\ \\

    \theutterance \stepcounter{utterance}  
    & & & \multicolumn{2}{p{0.3\linewidth}}{
        \cellcolor[rgb]{0.9,0.9,0.9}{
            \makecell[{{p{\linewidth}}}]{
                \texttt{\tiny{[GM$|$GM]}}
                \texttt{Ceza: (1,5) boş değil; "─" içeriyor.} \\
\\ 
\texttt{Toplamda 12 cezanın 2 kadarını biriktirdiniz.Lütfen tekrar deneyin!} \\
            }
        }
    }
    & & \\ \\

    \theutterance \stepcounter{utterance}  
    & & \multicolumn{4}{p{0.6\linewidth}}{
        \cellcolor[rgb]{0.9,0.9,0.9}{
            \makecell[{{p{\linewidth}}}]{
                \texttt{\tiny{[GM$\rangle$P2]}}
                \texttt{Ceza: (1,5) boş değil; "─" içeriyor.} \\
\\ 
\texttt{Toplamda 12 cezanın 2 kadarını biriktirdiniz.Lütfen tekrar deneyin!} \\
            }
        }
    }
    & \\ \\

    \theutterance \stepcounter{utterance}  
    & & & \multicolumn{4}{p{0.6\linewidth}}{
        \cellcolor[rgb]{0.9,0.9,0.9}{
            \makecell[{{p{\linewidth}}}]{
                \texttt{\tiny{[GM$\langle$P2]}}
                \texttt{HAMLE: W, (1, 6)} \\
            }
        }
    }
     \\ \\

    \theutterance \stepcounter{utterance}  
    & & & \multicolumn{2}{p{0.3\linewidth}}{
        \cellcolor[rgb]{0.9,0.9,0.9}{
            \makecell[{{p{\linewidth}}}]{
                \texttt{\tiny{[GM$|$GM]}}
                \texttt{HAMLE: W, (1, 6)} \\
            }
        }
    }
    & & \\ \\

    \theutterance \stepcounter{utterance}  
    & & & \multicolumn{2}{p{0.3\linewidth}}{
        \cellcolor[rgb]{0.9,0.9,0.9}{
            \makecell[{{p{\linewidth}}}]{
                \texttt{\tiny{[GM$|$GM]}}
                \texttt{Ceza: (1,6) boş değil; "H" içeriyor.} \\
\\ 
\texttt{Toplamda 12 cezanın 3 kadarını biriktirdiniz.Lütfen tekrar deneyin!} \\
            }
        }
    }
    & & \\ \\

    \theutterance \stepcounter{utterance}  
    & & \multicolumn{4}{p{0.6\linewidth}}{
        \cellcolor[rgb]{0.9,0.9,0.9}{
            \makecell[{{p{\linewidth}}}]{
                \texttt{\tiny{[GM$\rangle$P2]}}
                \texttt{Ceza: (1,6) boş değil; "H" içeriyor.} \\
\\ 
\texttt{Toplamda 12 cezanın 3 kadarını biriktirdiniz.Lütfen tekrar deneyin!} \\
            }
        }
    }
    & \\ \\

    \theutterance \stepcounter{utterance}  
    & & & \multicolumn{4}{p{0.6\linewidth}}{
        \cellcolor[rgb]{0.9,0.9,0.9}{
            \makecell[{{p{\linewidth}}}]{
                \texttt{\tiny{[GM$\langle$P2]}}
                \texttt{HAMLE: W, (2, 4)} \\
            }
        }
    }
     \\ \\

    \theutterance \stepcounter{utterance}  
    & & & \multicolumn{2}{p{0.3\linewidth}}{
        \cellcolor[rgb]{0.9,0.9,0.9}{
            \makecell[{{p{\linewidth}}}]{
                \texttt{\tiny{[GM$|$GM]}}
                \texttt{HAMLE: W, (2, 4)} \\
            }
        }
    }
    & & \\ \\

    \theutterance \stepcounter{utterance}  
    & & & \multicolumn{2}{p{0.3\linewidth}}{
        \cellcolor[rgb]{0.9,0.9,0.9}{
            \makecell[{{p{\linewidth}}}]{
                \texttt{\tiny{[GM$|$GM]}}
                \texttt{W moved to (2, 4)} \\
            }
        }
    }
    & & \\ \\

    \theutterance \stepcounter{utterance}  
    & & \multicolumn{4}{p{0.6\linewidth}}{
        \cellcolor[rgb]{0.9,0.9,0.9}{
            \makecell[{{p{\linewidth}}}]{
                \texttt{\tiny{[P1$\langle$GM]}}
                \texttt{T’ başarıyla (4,1) koordinatlarına taşındı. Güncellenmiş ızgaran şu şekilde:} \\
\texttt{```} \\
\texttt{ 1234567} \\
\texttt{╔════╤══╗} \\
\texttt{║H◌◌T│◌◌║ 1} \\
\texttt{║◌W◌◌│◌◌║ 2} \\
\texttt{║◌C◌◌│◌◌║ 3} \\
\texttt{║◌◌◌◌│I◌║ 4} \\
\texttt{╟──┐◌│◌◌║ 5} \\
\texttt{║◌◌├─┼──╢ 6} \\
\texttt{║◌◌│◌│◌◌║ 7} \\
\texttt{╚══╧═╧══╝} \\
\texttt{```} \\
\texttt{Şu anda en fazla 20 turdan 3. turu oynuyorsunuz.} \\
\texttt{Toplamda 12 cezanın 3 kadarını biriktirdiniz.Diğer oyuncu kendi ızgarasında bir nesneyi hareket ettirdi.} \\
\\ 
\\ 
\texttt{Bir sonraki komutun nedir?} \\
            }
        }
    }
    & \\ \\

    \theutterance \stepcounter{utterance}  
    & \multicolumn{4}{p{0.6\linewidth}}{
        \cellcolor[rgb]{0.9,0.9,0.9}{
            \makecell[{{p{\linewidth}}}]{
                \texttt{\tiny{[P1$\rangle$GM]}}
                \texttt{HAMLE: W, (1, 5)} \\
            }
        }
    }
    & & \\ \\

    \theutterance \stepcounter{utterance}  
    & & & \multicolumn{2}{p{0.3\linewidth}}{
        \cellcolor[rgb]{0.9,0.9,0.9}{
            \makecell[{{p{\linewidth}}}]{
                \texttt{\tiny{[GM$|$GM]}}
                \texttt{HAMLE: W, (1, 5)} \\
            }
        }
    }
    & & \\ \\

    \theutterance \stepcounter{utterance}  
    & & & \multicolumn{2}{p{0.3\linewidth}}{
        \cellcolor[rgb]{0.9,0.9,0.9}{
            \makecell[{{p{\linewidth}}}]{
                \texttt{\tiny{[GM$|$GM]}}
                \texttt{Ceza: (1,5) boş değil; "─" içeriyor.} \\
\\ 
\texttt{Toplamda 12 cezanın 4 kadarını biriktirdiniz.Lütfen tekrar deneyin!} \\
            }
        }
    }
    & & \\ \\

    \theutterance \stepcounter{utterance}  
    & & \multicolumn{4}{p{0.6\linewidth}}{
        \cellcolor[rgb]{0.9,0.9,0.9}{
            \makecell[{{p{\linewidth}}}]{
                \texttt{\tiny{[P1$\langle$GM]}}
                \texttt{Ceza: (1,5) boş değil; "─" içeriyor.} \\
\\ 
\texttt{Toplamda 12 cezanın 4 kadarını biriktirdiniz.Lütfen tekrar deneyin!} \\
            }
        }
    }
    & \\ \\

    \theutterance \stepcounter{utterance}  
    & \multicolumn{4}{p{0.6\linewidth}}{
        \cellcolor[rgb]{0.9,0.9,0.9}{
            \makecell[{{p{\linewidth}}}]{
                \texttt{\tiny{[P1$\rangle$GM]}}
                \texttt{HAMLE: W, (2, 5)} \\
            }
        }
    }
    & & \\ \\

    \theutterance \stepcounter{utterance}  
    & & & \multicolumn{2}{p{0.3\linewidth}}{
        \cellcolor[rgb]{0.9,0.9,0.9}{
            \makecell[{{p{\linewidth}}}]{
                \texttt{\tiny{[GM$|$GM]}}
                \texttt{HAMLE: W, (2, 5)} \\
            }
        }
    }
    & & \\ \\

    \theutterance \stepcounter{utterance}  
    & & & \multicolumn{2}{p{0.3\linewidth}}{
        \cellcolor[rgb]{0.9,0.9,0.9}{
            \makecell[{{p{\linewidth}}}]{
                \texttt{\tiny{[GM$|$GM]}}
                \texttt{Ceza: (2,5) boş değil; "─" içeriyor.} \\
\\ 
\texttt{Toplamda 12 cezanın 5 kadarını biriktirdiniz.Lütfen tekrar deneyin!} \\
            }
        }
    }
    & & \\ \\

    \theutterance \stepcounter{utterance}  
    & & \multicolumn{4}{p{0.6\linewidth}}{
        \cellcolor[rgb]{0.9,0.9,0.9}{
            \makecell[{{p{\linewidth}}}]{
                \texttt{\tiny{[P1$\langle$GM]}}
                \texttt{Ceza: (2,5) boş değil; "─" içeriyor.} \\
\\ 
\texttt{Toplamda 12 cezanın 5 kadarını biriktirdiniz.Lütfen tekrar deneyin!} \\
            }
        }
    }
    & \\ \\

    \theutterance \stepcounter{utterance}  
    & \multicolumn{4}{p{0.6\linewidth}}{
        \cellcolor[rgb]{0.9,0.9,0.9}{
            \makecell[{{p{\linewidth}}}]{
                \texttt{\tiny{[P1$\rangle$GM]}}
                \texttt{HAMLE: W, (3, 5)} \\
            }
        }
    }
    & & \\ \\

    \theutterance \stepcounter{utterance}  
    & & & \multicolumn{2}{p{0.3\linewidth}}{
        \cellcolor[rgb]{0.9,0.9,0.9}{
            \makecell[{{p{\linewidth}}}]{
                \texttt{\tiny{[GM$|$GM]}}
                \texttt{HAMLE: W, (3, 5)} \\
            }
        }
    }
    & & \\ \\

    \theutterance \stepcounter{utterance}  
    & & & \multicolumn{2}{p{0.3\linewidth}}{
        \cellcolor[rgb]{0.9,0.9,0.9}{
            \makecell[{{p{\linewidth}}}]{
                \texttt{\tiny{[GM$|$GM]}}
                \texttt{Ceza: (3,5) boş değil; "┐" içeriyor.} \\
\\ 
\texttt{Toplamda 12 cezanın 6 kadarını biriktirdiniz.Lütfen tekrar deneyin!} \\
            }
        }
    }
    & & \\ \\

    \theutterance \stepcounter{utterance}  
    & & \multicolumn{4}{p{0.6\linewidth}}{
        \cellcolor[rgb]{0.9,0.9,0.9}{
            \makecell[{{p{\linewidth}}}]{
                \texttt{\tiny{[P1$\langle$GM]}}
                \texttt{Ceza: (3,5) boş değil; "┐" içeriyor.} \\
\\ 
\texttt{Toplamda 12 cezanın 6 kadarını biriktirdiniz.Lütfen tekrar deneyin!} \\
            }
        }
    }
    & \\ \\

    \theutterance \stepcounter{utterance}  
    & \multicolumn{4}{p{0.6\linewidth}}{
        \cellcolor[rgb]{0.9,0.9,0.9}{
            \makecell[{{p{\linewidth}}}]{
                \texttt{\tiny{[P1$\rangle$GM]}}
                \texttt{HAMLE: W, (4, 5)} \\
            }
        }
    }
    & & \\ \\

    \theutterance \stepcounter{utterance}  
    & & & \multicolumn{2}{p{0.3\linewidth}}{
        \cellcolor[rgb]{0.9,0.9,0.9}{
            \makecell[{{p{\linewidth}}}]{
                \texttt{\tiny{[GM$|$GM]}}
                \texttt{HAMLE: W, (4, 5)} \\
            }
        }
    }
    & & \\ \\

    \theutterance \stepcounter{utterance}  
    & & & \multicolumn{2}{p{0.3\linewidth}}{
        \cellcolor[rgb]{0.9,0.9,0.9}{
            \makecell[{{p{\linewidth}}}]{
                \texttt{\tiny{[GM$|$GM]}}
                \texttt{W moved to (4, 5)} \\
            }
        }
    }
    & & \\ \\

    \theutterance \stepcounter{utterance}  
    & & \multicolumn{4}{p{0.6\linewidth}}{
        \cellcolor[rgb]{0.9,0.9,0.9}{
            \makecell[{{p{\linewidth}}}]{
                \texttt{\tiny{[GM$\rangle$P2]}}
                \texttt{W’ başarıyla (2,4) koordinatlarına taşındı. Güncellenmiş ızgaran şu şekilde:} \\
\texttt{```} \\
\texttt{ 1234567} \\
\texttt{╔════╤══╗} \\
\texttt{║◌◌◌T│◌◌║ 1} \\
\texttt{║◌◌◌◌│◌◌║ 2} \\
\texttt{║◌◌◌◌│◌◌║ 3} \\
\texttt{║◌W◌◌│◌◌║ 4} \\
\texttt{╟──┐◌│◌◌║ 5} \\
\texttt{║H◌├─┼──╢ 6} \\
\texttt{║I◌│◌│◌C║ 7} \\
\texttt{╚══╧═╧══╝} \\
\texttt{```} \\
\texttt{Şu anda en fazla 20 turdan 4. turu oynuyorsunuz.} \\
\texttt{Toplamda 12 cezanın 6 kadarını biriktirdiniz.Diğer oyuncu kendi ızgarasında bir nesneyi hareket ettirdi.} \\
\\ 
\\ 
\texttt{Bir sonraki komutun nedir?} \\
            }
        }
    }
    & \\ \\

    \theutterance \stepcounter{utterance}  
    & & & \multicolumn{4}{p{0.6\linewidth}}{
        \cellcolor[rgb]{0.9,0.9,0.9}{
            \makecell[{{p{\linewidth}}}]{
                \texttt{\tiny{[GM$\langle$P2]}}
                \texttt{SÖYLE: W'yi (2,4) konumuna taşıdım. Şimdi H'yi (1,5) konumuna taşıyorum.} \\
            }
        }
    }
     \\ \\

    \theutterance \stepcounter{utterance}  
    & & & \multicolumn{2}{p{0.3\linewidth}}{
        \cellcolor[rgb]{0.9,0.9,0.9}{
            \makecell[{{p{\linewidth}}}]{
                \texttt{\tiny{[GM$|$GM]}}
                \texttt{SÖYLE: W'yi (2,4) konumuna taşıdım. Şimdi H'yi (1,5) konumuna taşıyorum.} \\
            }
        }
    }
    & & \\ \\

    \theutterance \stepcounter{utterance}  
    & & \multicolumn{4}{p{0.6\linewidth}}{
        \cellcolor[rgb]{0.9,0.9,0.9}{
            \makecell[{{p{\linewidth}}}]{
                \texttt{\tiny{[P1$\langle$GM]}}
                \texttt{W’ başarıyla (4,5) koordinatlarına taşındı. Güncellenmiş ızgaran şu şekilde:} \\
\texttt{```} \\
\texttt{ 1234567} \\
\texttt{╔════╤══╗} \\
\texttt{║H◌◌T│◌◌║ 1} \\
\texttt{║◌◌◌◌│◌◌║ 2} \\
\texttt{║◌C◌◌│◌◌║ 3} \\
\texttt{║◌◌◌◌│I◌║ 4} \\
\texttt{╟──┐W│◌◌║ 5} \\
\texttt{║◌◌├─┼──╢ 6} \\
\texttt{║◌◌│◌│◌◌║ 7} \\
\texttt{╚══╧═╧══╝} \\
\texttt{```} \\
\texttt{Şu anda en fazla 20 turdan 4. turu oynuyorsunuz.} \\
\texttt{Toplamda 12 cezanın 6 kadarını biriktirdiniz.Diğer oyuncu şu mesajı gönderdi:} \\
\\ 
\\ 
\texttt{"W'yi (2,4) konumuna taşıdım. Şimdi H'yi (1,5) konumuna taşıyorum."} \\
\\ 
\\ 
\texttt{Bir sonraki komutun nedir?} \\
            }
        }
    }
    & \\ \\

    \theutterance \stepcounter{utterance}  
    & \multicolumn{4}{p{0.6\linewidth}}{
        \cellcolor[rgb]{0.9,0.9,0.9}{
            \makecell[{{p{\linewidth}}}]{
                \texttt{\tiny{[P1$\rangle$GM]}}
                \texttt{SÖYLE: Dikkat, benim tahtamda 5. satır (y=5) tamamen duvarlarla dolu (x=3,4,5,6,7). H için (1,5) mümkün değil. H'yi (1,4)'e taşıyorum, o hücre boş. (1,4) senin tahtanda da boş mu? Uygunysa H hedefi (1,4) olsun.} \\
            }
        }
    }
    & & \\ \\

    \theutterance \stepcounter{utterance}  
    & & & \multicolumn{2}{p{0.3\linewidth}}{
        \cellcolor[rgb]{0.9,0.9,0.9}{
            \makecell[{{p{\linewidth}}}]{
                \texttt{\tiny{[GM$|$GM]}}
                \texttt{SÖYLE: Dikkat, benim tahtamda 5. satır (y=5) tamamen duvarlarla dolu (x=3,4,5,6,7). H için (1,5) mümkün değil. H'yi (1,4)'e taşıyorum, o hücre boş. (1,4) senin tahtanda da boş mu? Uygunysa H hedefi (1,4) olsun.} \\
            }
        }
    }
    & & \\ \\

    \theutterance \stepcounter{utterance}  
    & & \multicolumn{4}{p{0.6\linewidth}}{
        \cellcolor[rgb]{0.9,0.9,0.9}{
            \makecell[{{p{\linewidth}}}]{
                \texttt{\tiny{[GM$\rangle$P2]}}
                \texttt{Mesajın diğer oyuncuya iletildi.} \\
\\ 
\\ 
\texttt{Şu anda en fazla 20 turdan 5. turu oynuyorsunuz.} \\
\texttt{Toplamda 12 cezanın 6 kadarını biriktirdiniz.Diğer oyuncu şu mesajı gönderdi:} \\
\\ 
\\ 
\texttt{"Dikkat, benim tahtamda 5. satır (y=5) tamamen duvarlarla dolu (x=3,4,5,6,7). H için (1,5) mümkün değil. H'yi (1,4)'e taşıyorum, o hücre boş. (1,4) senin tahtanda da boş mu? Uygunysa H hedefi (1,4) olsun."} \\
\\ 
\\ 
\texttt{Bir sonraki komutun nedir?} \\
            }
        }
    }
    & \\ \\

    \theutterance \stepcounter{utterance}  
    & & & \multicolumn{4}{p{0.6\linewidth}}{
        \cellcolor[rgb]{0.9,0.9,0.9}{
            \makecell[{{p{\linewidth}}}]{
                \texttt{\tiny{[GM$\langle$P2]}}
                \texttt{HAMLE: H, (1, 4)} \\
            }
        }
    }
     \\ \\

    \theutterance \stepcounter{utterance}  
    & & & \multicolumn{2}{p{0.3\linewidth}}{
        \cellcolor[rgb]{0.9,0.9,0.9}{
            \makecell[{{p{\linewidth}}}]{
                \texttt{\tiny{[GM$|$GM]}}
                \texttt{HAMLE: H, (1, 4)} \\
            }
        }
    }
    & & \\ \\

    \theutterance \stepcounter{utterance}  
    & & & \multicolumn{2}{p{0.3\linewidth}}{
        \cellcolor[rgb]{0.9,0.9,0.9}{
            \makecell[{{p{\linewidth}}}]{
                \texttt{\tiny{[GM$|$GM]}}
                \texttt{H moved to (1, 4)} \\
            }
        }
    }
    & & \\ \\

    \theutterance \stepcounter{utterance}  
    & & \multicolumn{4}{p{0.6\linewidth}}{
        \cellcolor[rgb]{0.9,0.9,0.9}{
            \makecell[{{p{\linewidth}}}]{
                \texttt{\tiny{[P1$\langle$GM]}}
                \texttt{Mesajın diğer oyuncuya iletildi.} \\
\\ 
\\ 
\texttt{Şu anda en fazla 20 turdan 5. turu oynuyorsunuz.} \\
\texttt{Toplamda 12 cezanın 6 kadarını biriktirdiniz.Diğer oyuncu kendi ızgarasında bir nesneyi hareket ettirdi.} \\
\\ 
\\ 
\texttt{Bir sonraki komutun nedir?} \\
            }
        }
    }
    & \\ \\

    \theutterance \stepcounter{utterance}  
    & \multicolumn{4}{p{0.6\linewidth}}{
        \cellcolor[rgb]{0.9,0.9,0.9}{
            \makecell[{{p{\linewidth}}}]{
                \texttt{\tiny{[P1$\rangle$GM]}}
                \texttt{HAMLE: H, (1, 4)} \\
            }
        }
    }
    & & \\ \\

    \theutterance \stepcounter{utterance}  
    & & & \multicolumn{2}{p{0.3\linewidth}}{
        \cellcolor[rgb]{0.9,0.9,0.9}{
            \makecell[{{p{\linewidth}}}]{
                \texttt{\tiny{[GM$|$GM]}}
                \texttt{HAMLE: H, (1, 4)} \\
            }
        }
    }
    & & \\ \\

    \theutterance \stepcounter{utterance}  
    & & & \multicolumn{2}{p{0.3\linewidth}}{
        \cellcolor[rgb]{0.9,0.9,0.9}{
            \makecell[{{p{\linewidth}}}]{
                \texttt{\tiny{[GM$|$GM]}}
                \texttt{H moved to (1, 4)} \\
            }
        }
    }
    & & \\ \\

    \theutterance \stepcounter{utterance}  
    & & \multicolumn{4}{p{0.6\linewidth}}{
        \cellcolor[rgb]{0.9,0.9,0.9}{
            \makecell[{{p{\linewidth}}}]{
                \texttt{\tiny{[GM$\rangle$P2]}}
                \texttt{H’ başarıyla (1,4) koordinatlarına taşındı. Güncellenmiş ızgaran şu şekilde:} \\
\texttt{```} \\
\texttt{ 1234567} \\
\texttt{╔════╤══╗} \\
\texttt{║◌◌◌T│◌◌║ 1} \\
\texttt{║◌◌◌◌│◌◌║ 2} \\
\texttt{║◌◌◌◌│◌◌║ 3} \\
\texttt{║HW◌◌│◌◌║ 4} \\
\texttt{╟──┐◌│◌◌║ 5} \\
\texttt{║◌◌├─┼──╢ 6} \\
\texttt{║I◌│◌│◌C║ 7} \\
\texttt{╚══╧═╧══╝} \\
\texttt{```} \\
\texttt{Şu anda en fazla 20 turdan 6. turu oynuyorsunuz.} \\
\texttt{Toplamda 12 cezanın 6 kadarını biriktirdiniz.Diğer oyuncu kendi ızgarasında bir nesneyi hareket ettirdi.} \\
\\ 
\\ 
\texttt{Bir sonraki komutun nedir?} \\
            }
        }
    }
    & \\ \\

    \theutterance \stepcounter{utterance}  
    & & & \multicolumn{4}{p{0.6\linewidth}}{
        \cellcolor[rgb]{0.9,0.9,0.9}{
            \makecell[{{p{\linewidth}}}]{
                \texttt{\tiny{[GM$\langle$P2]}}
                \texttt{SÖYLE: H'yi (1,4) konumuna taşıdım. H hedefi (1,4) olarak güncellendi. Şimdi I'yi (7,6) konumuna taşımaya çalışıyorum.} \\
            }
        }
    }
     \\ \\

    \theutterance \stepcounter{utterance}  
    & & & \multicolumn{2}{p{0.3\linewidth}}{
        \cellcolor[rgb]{0.9,0.9,0.9}{
            \makecell[{{p{\linewidth}}}]{
                \texttt{\tiny{[GM$|$GM]}}
                \texttt{SÖYLE: H'yi (1,4) konumuna taşıdım. H hedefi (1,4) olarak güncellendi. Şimdi I'yi (7,6) konumuna taşımaya çalışıyorum.} \\
            }
        }
    }
    & & \\ \\

    \theutterance \stepcounter{utterance}  
    & & \multicolumn{4}{p{0.6\linewidth}}{
        \cellcolor[rgb]{0.9,0.9,0.9}{
            \makecell[{{p{\linewidth}}}]{
                \texttt{\tiny{[P1$\langle$GM]}}
                \texttt{H’ başarıyla (1,4) koordinatlarına taşındı. Güncellenmiş ızgaran şu şekilde:} \\
\texttt{```} \\
\texttt{ 1234567} \\
\texttt{╔════╤══╗} \\
\texttt{║◌◌◌T│◌◌║ 1} \\
\texttt{║◌◌◌◌│◌◌║ 2} \\
\texttt{║◌C◌◌│◌◌║ 3} \\
\texttt{║H◌◌◌│I◌║ 4} \\
\texttt{╟──┐W│◌◌║ 5} \\
\texttt{║◌◌├─┼──╢ 6} \\
\texttt{║◌◌│◌│◌◌║ 7} \\
\texttt{╚══╧═╧══╝} \\
\texttt{```} \\
\texttt{Şu anda en fazla 20 turdan 6. turu oynuyorsunuz.} \\
\texttt{Toplamda 12 cezanın 6 kadarını biriktirdiniz.Diğer oyuncu şu mesajı gönderdi:} \\
\\ 
\\ 
\texttt{"H'yi (1,4) konumuna taşıdım. H hedefi (1,4) olarak güncellendi. Şimdi I'yi (7,6) konumuna taşımaya çalışıyorum."} \\
\\ 
\\ 
\texttt{Bir sonraki komutun nedir?} \\
            }
        }
    }
    & \\ \\

    \theutterance \stepcounter{utterance}  
    & \multicolumn{4}{p{0.6\linewidth}}{
        \cellcolor[rgb]{0.9,0.9,0.9}{
            \makecell[{{p{\linewidth}}}]{
                \texttt{\tiny{[P1$\rangle$GM]}}
                \texttt{HAMLE: I, (7, 6)} \\
            }
        }
    }
    & & \\ \\

    \theutterance \stepcounter{utterance}  
    & & & \multicolumn{2}{p{0.3\linewidth}}{
        \cellcolor[rgb]{0.9,0.9,0.9}{
            \makecell[{{p{\linewidth}}}]{
                \texttt{\tiny{[GM$|$GM]}}
                \texttt{HAMLE: I, (7, 6)} \\
            }
        }
    }
    & & \\ \\

    \theutterance \stepcounter{utterance}  
    & & & \multicolumn{2}{p{0.3\linewidth}}{
        \cellcolor[rgb]{0.9,0.9,0.9}{
            \makecell[{{p{\linewidth}}}]{
                \texttt{\tiny{[GM$|$GM]}}
                \texttt{Ceza: (7,6) boş değil; "─" içeriyor.} \\
\\ 
\texttt{Toplamda 12 cezanın 7 kadarını biriktirdiniz.Lütfen tekrar deneyin!} \\
            }
        }
    }
    & & \\ \\

    \theutterance \stepcounter{utterance}  
    & & \multicolumn{4}{p{0.6\linewidth}}{
        \cellcolor[rgb]{0.9,0.9,0.9}{
            \makecell[{{p{\linewidth}}}]{
                \texttt{\tiny{[P1$\langle$GM]}}
                \texttt{Ceza: (7,6) boş değil; "─" içeriyor.} \\
\\ 
\texttt{Toplamda 12 cezanın 7 kadarını biriktirdiniz.Lütfen tekrar deneyin!} \\
            }
        }
    }
    & \\ \\

    \theutterance \stepcounter{utterance}  
    & \multicolumn{4}{p{0.6\linewidth}}{
        \cellcolor[rgb]{0.9,0.9,0.9}{
            \makecell[{{p{\linewidth}}}]{
                \texttt{\tiny{[P1$\rangle$GM]}}
                \texttt{HAMLE: I, (7, 5)} \\
            }
        }
    }
    & & \\ \\

    \theutterance \stepcounter{utterance}  
    & & & \multicolumn{2}{p{0.3\linewidth}}{
        \cellcolor[rgb]{0.9,0.9,0.9}{
            \makecell[{{p{\linewidth}}}]{
                \texttt{\tiny{[GM$|$GM]}}
                \texttt{HAMLE: I, (7, 5)} \\
            }
        }
    }
    & & \\ \\

    \theutterance \stepcounter{utterance}  
    & & & \multicolumn{2}{p{0.3\linewidth}}{
        \cellcolor[rgb]{0.9,0.9,0.9}{
            \makecell[{{p{\linewidth}}}]{
                \texttt{\tiny{[GM$|$GM]}}
                \texttt{I moved to (7, 5)} \\
            }
        }
    }
    & & \\ \\

    \theutterance \stepcounter{utterance}  
    & & \multicolumn{4}{p{0.6\linewidth}}{
        \cellcolor[rgb]{0.9,0.9,0.9}{
            \makecell[{{p{\linewidth}}}]{
                \texttt{\tiny{[GM$\rangle$P2]}}
                \texttt{Mesajın diğer oyuncuya iletildi.} \\
\\ 
\\ 
\texttt{Şu anda en fazla 20 turdan 7. turu oynuyorsunuz.} \\
\texttt{Toplamda 12 cezanın 7 kadarını biriktirdiniz.Diğer oyuncu kendi ızgarasında bir nesneyi hareket ettirdi.} \\
\\ 
\\ 
\texttt{Bir sonraki komutun nedir?} \\
            }
        }
    }
    & \\ \\

    \theutterance \stepcounter{utterance}  
    & & & \multicolumn{4}{p{0.6\linewidth}}{
        \cellcolor[rgb]{0.9,0.9,0.9}{
            \makecell[{{p{\linewidth}}}]{
                \texttt{\tiny{[GM$\langle$P2]}}
                \texttt{HAMLE: I, (7, 6)} \\
            }
        }
    }
     \\ \\

    \theutterance \stepcounter{utterance}  
    & & & \multicolumn{2}{p{0.3\linewidth}}{
        \cellcolor[rgb]{0.9,0.9,0.9}{
            \makecell[{{p{\linewidth}}}]{
                \texttt{\tiny{[GM$|$GM]}}
                \texttt{HAMLE: I, (7, 6)} \\
            }
        }
    }
    & & \\ \\

    \theutterance \stepcounter{utterance}  
    & & & \multicolumn{2}{p{0.3\linewidth}}{
        \cellcolor[rgb]{0.9,0.9,0.9}{
            \makecell[{{p{\linewidth}}}]{
                \texttt{\tiny{[GM$|$GM]}}
                \texttt{Ceza: (7,6) boş değil; "─" içeriyor.} \\
\\ 
\texttt{Toplamda 12 cezanın 8 kadarını biriktirdiniz.Lütfen tekrar deneyin!} \\
            }
        }
    }
    & & \\ \\

    \theutterance \stepcounter{utterance}  
    & & \multicolumn{4}{p{0.6\linewidth}}{
        \cellcolor[rgb]{0.9,0.9,0.9}{
            \makecell[{{p{\linewidth}}}]{
                \texttt{\tiny{[GM$\rangle$P2]}}
                \texttt{Ceza: (7,6) boş değil; "─" içeriyor.} \\
\\ 
\texttt{Toplamda 12 cezanın 8 kadarını biriktirdiniz.Lütfen tekrar deneyin!} \\
            }
        }
    }
    & \\ \\

    \theutterance \stepcounter{utterance}  
    & & & \multicolumn{4}{p{0.6\linewidth}}{
        \cellcolor[rgb]{0.9,0.9,0.9}{
            \makecell[{{p{\linewidth}}}]{
                \texttt{\tiny{[GM$\langle$P2]}}
                \texttt{HAMLE: I, (6, 6)} \\
            }
        }
    }
     \\ \\

    \theutterance \stepcounter{utterance}  
    & & & \multicolumn{2}{p{0.3\linewidth}}{
        \cellcolor[rgb]{0.9,0.9,0.9}{
            \makecell[{{p{\linewidth}}}]{
                \texttt{\tiny{[GM$|$GM]}}
                \texttt{HAMLE: I, (6, 6)} \\
            }
        }
    }
    & & \\ \\

    \theutterance \stepcounter{utterance}  
    & & & \multicolumn{2}{p{0.3\linewidth}}{
        \cellcolor[rgb]{0.9,0.9,0.9}{
            \makecell[{{p{\linewidth}}}]{
                \texttt{\tiny{[GM$|$GM]}}
                \texttt{Ceza: (6,6) boş değil; "─" içeriyor.} \\
\\ 
\texttt{Toplamda 12 cezanın 9 kadarını biriktirdiniz.Lütfen tekrar deneyin!} \\
            }
        }
    }
    & & \\ \\

    \theutterance \stepcounter{utterance}  
    & & \multicolumn{4}{p{0.6\linewidth}}{
        \cellcolor[rgb]{0.9,0.9,0.9}{
            \makecell[{{p{\linewidth}}}]{
                \texttt{\tiny{[GM$\rangle$P2]}}
                \texttt{Ceza: (6,6) boş değil; "─" içeriyor.} \\
\\ 
\texttt{Toplamda 12 cezanın 9 kadarını biriktirdiniz.Lütfen tekrar deneyin!} \\
            }
        }
    }
    & \\ \\

    \theutterance \stepcounter{utterance}  
    & & & \multicolumn{4}{p{0.6\linewidth}}{
        \cellcolor[rgb]{0.9,0.9,0.9}{
            \makecell[{{p{\linewidth}}}]{
                \texttt{\tiny{[GM$\langle$P2]}}
                \texttt{HAMLE: I, (5, 6)} \\
            }
        }
    }
     \\ \\

    \theutterance \stepcounter{utterance}  
    & & & \multicolumn{2}{p{0.3\linewidth}}{
        \cellcolor[rgb]{0.9,0.9,0.9}{
            \makecell[{{p{\linewidth}}}]{
                \texttt{\tiny{[GM$|$GM]}}
                \texttt{HAMLE: I, (5, 6)} \\
            }
        }
    }
    & & \\ \\

    \theutterance \stepcounter{utterance}  
    & & & \multicolumn{2}{p{0.3\linewidth}}{
        \cellcolor[rgb]{0.9,0.9,0.9}{
            \makecell[{{p{\linewidth}}}]{
                \texttt{\tiny{[GM$|$GM]}}
                \texttt{Ceza: (5,6) boş değil; "┼" içeriyor.} \\
\\ 
\texttt{Toplamda 12 cezanın 10 kadarını biriktirdiniz.Lütfen tekrar deneyin!} \\
            }
        }
    }
    & & \\ \\

    \theutterance \stepcounter{utterance}  
    & & \multicolumn{4}{p{0.6\linewidth}}{
        \cellcolor[rgb]{0.9,0.9,0.9}{
            \makecell[{{p{\linewidth}}}]{
                \texttt{\tiny{[GM$\rangle$P2]}}
                \texttt{Ceza: (5,6) boş değil; "┼" içeriyor.} \\
\\ 
\texttt{Toplamda 12 cezanın 10 kadarını biriktirdiniz.Lütfen tekrar deneyin!} \\
            }
        }
    }
    & \\ \\

    \theutterance \stepcounter{utterance}  
    & & & \multicolumn{4}{p{0.6\linewidth}}{
        \cellcolor[rgb]{0.9,0.9,0.9}{
            \makecell[{{p{\linewidth}}}]{
                \texttt{\tiny{[GM$\langle$P2]}}
                \texttt{HAMLE: I, (4, 6)} \\
            }
        }
    }
     \\ \\

    \theutterance \stepcounter{utterance}  
    & & & \multicolumn{2}{p{0.3\linewidth}}{
        \cellcolor[rgb]{0.9,0.9,0.9}{
            \makecell[{{p{\linewidth}}}]{
                \texttt{\tiny{[GM$|$GM]}}
                \texttt{HAMLE: I, (4, 6)} \\
            }
        }
    }
    & & \\ \\

    \theutterance \stepcounter{utterance}  
    & & & \multicolumn{2}{p{0.3\linewidth}}{
        \cellcolor[rgb]{0.9,0.9,0.9}{
            \makecell[{{p{\linewidth}}}]{
                \texttt{\tiny{[GM$|$GM]}}
                \texttt{Ceza: (4,6) boş değil; "─" içeriyor.} \\
\\ 
\texttt{Toplamda 12 cezanın 11 kadarını biriktirdiniz.Lütfen tekrar deneyin!} \\
            }
        }
    }
    & & \\ \\

    \theutterance \stepcounter{utterance}  
    & & \multicolumn{4}{p{0.6\linewidth}}{
        \cellcolor[rgb]{0.9,0.9,0.9}{
            \makecell[{{p{\linewidth}}}]{
                \texttt{\tiny{[GM$\rangle$P2]}}
                \texttt{Ceza: (4,6) boş değil; "─" içeriyor.} \\
\\ 
\texttt{Toplamda 12 cezanın 11 kadarını biriktirdiniz.Lütfen tekrar deneyin!} \\
            }
        }
    }
    & \\ \\

    \theutterance \stepcounter{utterance}  
    & & & \multicolumn{4}{p{0.6\linewidth}}{
        \cellcolor[rgb]{0.9,0.9,0.9}{
            \makecell[{{p{\linewidth}}}]{
                \texttt{\tiny{[GM$\langle$P2]}}
                \texttt{HAMLE: I, (3, 6)} \\
            }
        }
    }
     \\ \\

    \theutterance \stepcounter{utterance}  
    & & & \multicolumn{2}{p{0.3\linewidth}}{
        \cellcolor[rgb]{0.9,0.9,0.9}{
            \makecell[{{p{\linewidth}}}]{
                \texttt{\tiny{[GM$|$GM]}}
                \texttt{HAMLE: I, (3, 6)} \\
            }
        }
    }
    & & \\ \\

    \theutterance \stepcounter{utterance}  
    & & & \multicolumn{2}{p{0.3\linewidth}}{
        \cellcolor[rgb]{0.9,0.9,0.9}{
            \makecell[{{p{\linewidth}}}]{
                \texttt{\tiny{[GM$|$GM]}}
                \texttt{Ceza: (3,6) boş değil; "├" içeriyor.} \\
\\ 
\texttt{Toplamda 12 cezanın 12 kadarını biriktirdiniz.Lütfen tekrar deneyin!} \\
            }
        }
    }
    & & \\ \\

    \theutterance \stepcounter{utterance}  
    & & \multicolumn{4}{p{0.6\linewidth}}{
        \cellcolor[rgb]{0.9,0.9,0.9}{
            \makecell[{{p{\linewidth}}}]{
                \texttt{\tiny{[GM$\rangle$P2]}}
                \texttt{Ceza: (3,6) boş değil; "├" içeriyor.} \\
\\ 
\texttt{Toplamda 12 cezanın 12 kadarını biriktirdiniz.Lütfen tekrar deneyin!} \\
            }
        }
    }
    & \\ \\

    \theutterance \stepcounter{utterance}  
    & & & \multicolumn{4}{p{0.6\linewidth}}{
        \cellcolor[rgb]{0.9,0.9,0.9}{
            \makecell[{{p{\linewidth}}}]{
                \texttt{\tiny{[GM$\langle$P2]}}
                \texttt{SÖYLE: Benim tahtamda 6. satır (y=6) x=3'ten x=7'ye kadar tamamen duvarlarla dolu (├─┼──). I'yi y=6 satırına koyamıyorum. I için yeni hedef önerisi: (2, 5) veya (2, 4) civarı nasıl? Sende I nerede duruyor?} \\
            }
        }
    }
     \\ \\

    \theutterance \stepcounter{utterance}  
    & & & \multicolumn{2}{p{0.3\linewidth}}{
        \cellcolor[rgb]{0.9,0.9,0.9}{
            \makecell[{{p{\linewidth}}}]{
                \texttt{\tiny{[GM$|$GM]}}
                \texttt{SÖYLE: Benim tahtamda 6. satır (y=6) x=3'ten x=7'ye kadar tamamen duvarlarla dolu (├─┼──). I'yi y=6 satırına koyamıyorum. I için yeni hedef önerisi: (2, 5) veya (2, 4) civarı nasıl? Sende I nerede duruyor?} \\
            }
        }
    }
    & & \\ \\

    \theutterance \stepcounter{utterance}  
    & & \multicolumn{4}{p{0.6\linewidth}}{
        \cellcolor[rgb]{0.9,0.9,0.9}{
            \makecell[{{p{\linewidth}}}]{
                \texttt{\tiny{[P1$\langle$GM]}}
                \texttt{I’ başarıyla (7,5) koordinatlarına taşındı. Güncellenmiş ızgaran şu şekilde:} \\
\texttt{```} \\
\texttt{ 1234567} \\
\texttt{╔════╤══╗} \\
\texttt{║◌◌◌T│◌◌║ 1} \\
\texttt{║◌◌◌◌│◌◌║ 2} \\
\texttt{║◌C◌◌│◌◌║ 3} \\
\texttt{║H◌◌◌│◌◌║ 4} \\
\texttt{╟──┐W│◌I║ 5} \\
\texttt{║◌◌├─┼──╢ 6} \\
\texttt{║◌◌│◌│◌◌║ 7} \\
\texttt{╚══╧═╧══╝} \\
\texttt{```} \\
\texttt{Şu anda en fazla 20 turdan 7. turu oynuyorsunuz.} \\
\texttt{Toplamda 12 cezanın 12 kadarını biriktirdiniz.Diğer oyuncu şu mesajı gönderdi:} \\
\\ 
\\ 
\texttt{"Benim tahtamda 6. satır (y=6) x=3'ten x=7'ye kadar tamamen duvarlarla dolu (├─┼──). I'yi y=6 satırına koyamıyorum. I için yeni hedef önerisi: (2, 5) veya (2, 4) civarı nasıl? Sende I nerede duruyor?"} \\
\\ 
\\ 
\texttt{Bir sonraki komutun nedir?} \\
            }
        }
    }
    & \\ \\

    \theutterance \stepcounter{utterance}  
    & \multicolumn{4}{p{0.6\linewidth}}{
        \cellcolor[rgb]{0.9,0.9,0.9}{
            \makecell[{{p{\linewidth}}}]{
                \texttt{\tiny{[P1$\rangle$GM]}}
                \texttt{SÖYLE: Benim tahtamda I şu an (7,5) konumunda. (2,5) benim tahtamda duvar (─). (2,4) boş. C için (7,7) hedefi boş. Yeni öneri: H(1,4), T(4,1), W(4,5), I(2,4), C(7,7). Bu koordinatlar benim tahtamda boş. Senin tahtanda nasıl?} \\
            }
        }
    }
    & & \\ \\

    \theutterance \stepcounter{utterance}  
    & & & \multicolumn{2}{p{0.3\linewidth}}{
        \cellcolor[rgb]{0.9,0.9,0.9}{
            \makecell[{{p{\linewidth}}}]{
                \texttt{\tiny{[GM$|$GM]}}
                \texttt{SÖYLE: Benim tahtamda I şu an (7,5) konumunda. (2,5) benim tahtamda duvar (─). (2,4) boş. C için (7,7) hedefi boş. Yeni öneri: H(1,4), T(4,1), W(4,5), I(2,4), C(7,7). Bu koordinatlar benim tahtamda boş. Senin tahtanda nasıl?} \\
            }
        }
    }
    & & \\ \\

    \theutterance \stepcounter{utterance}  
    & & \multicolumn{4}{p{0.6\linewidth}}{
        \cellcolor[rgb]{0.9,0.9,0.9}{
            \makecell[{{p{\linewidth}}}]{
                \texttt{\tiny{[GM$\rangle$P2]}}
                \texttt{Mesajın diğer oyuncuya iletildi.} \\
\\ 
\\ 
\texttt{Şu anda en fazla 20 turdan 8. turu oynuyorsunuz.} \\
\texttt{Toplamda 12 cezanın 12 kadarını biriktirdiniz.Diğer oyuncu şu mesajı gönderdi:} \\
\\ 
\\ 
\texttt{"Benim tahtamda I şu an (7,5) konumunda. (2,5) benim tahtamda duvar (─). (2,4) boş. C için (7,7) hedefi boş. Yeni öneri: H(1,4), T(4,1), W(4,5), I(2,4), C(7,7). Bu koordinatlar benim tahtamda boş. Senin tahtanda nasıl?"} \\
\\ 
\\ 
\texttt{Bir sonraki komutun nedir?} \\
            }
        }
    }
    & \\ \\

    \theutterance \stepcounter{utterance}  
    & & & \multicolumn{4}{p{0.6\linewidth}}{
        \cellcolor[rgb]{0.9,0.9,0.9}{
            \makecell[{{p{\linewidth}}}]{
                \texttt{\tiny{[GM$\langle$P2]}}
                \texttt{SÖYLE: Yeni hedefler uygun: H(1,4) {[}yerinde{]}, T(4,1) {[}yerinde{]}, W(4,5), I(2,4), C(7,7). Benim tahtamda (4,5), (2,4), (7,7) boş. W'yi (4,5)'e, I'yi (2,4)'e, C'yi (7,7)'ye taşıyorum. Önce W: (4,5).} \\
            }
        }
    }
     \\ \\

    \theutterance \stepcounter{utterance}  
    & & & \multicolumn{2}{p{0.3\linewidth}}{
        \cellcolor[rgb]{0.9,0.9,0.9}{
            \makecell[{{p{\linewidth}}}]{
                \texttt{\tiny{[GM$|$GM]}}
                \texttt{SÖYLE: Yeni hedefler uygun: H(1,4) {[}yerinde{]}, T(4,1) {[}yerinde{]}, W(4,5), I(2,4), C(7,7). Benim tahtamda (4,5), (2,4), (7,7) boş. W'yi (4,5)'e, I'yi (2,4)'e, C'yi (7,7)'ye taşıyorum. Önce W: (4,5).} \\
            }
        }
    }
    & & \\ \\

    \theutterance \stepcounter{utterance}  
    & & \multicolumn{4}{p{0.6\linewidth}}{
        \cellcolor[rgb]{0.9,0.9,0.9}{
            \makecell[{{p{\linewidth}}}]{
                \texttt{\tiny{[P1$\langle$GM]}}
                \texttt{Mesajın diğer oyuncuya iletildi.} \\
\\ 
\\ 
\texttt{Şu anda en fazla 20 turdan 8. turu oynuyorsunuz.} \\
\texttt{Toplamda 12 cezanın 12 kadarını biriktirdiniz.Diğer oyuncu şu mesajı gönderdi:} \\
\\ 
\\ 
\texttt{"Yeni hedefler uygun: H(1,4) {[}yerinde{]}, T(4,1) {[}yerinde{]}, W(4,5), I(2,4), C(7,7). Benim tahtamda (4,5), (2,4), (7,7) boş. W'yi (4,5)'e, I'yi (2,4)'e, C'yi (7,7)'ye taşıyorum. Önce W: (4,5)."} \\
\\ 
\\ 
\texttt{Bir sonraki komutun nedir?} \\
            }
        }
    }
    & \\ \\

    \theutterance \stepcounter{utterance}  
    & \multicolumn{4}{p{0.6\linewidth}}{
        \cellcolor[rgb]{0.9,0.9,0.9}{
            \makecell[{{p{\linewidth}}}]{
                \texttt{\tiny{[P1$\rangle$GM]}}
                \texttt{HAMLE: C, (7, 7)} \\
            }
        }
    }
    & & \\ \\

    \theutterance \stepcounter{utterance}  
    & & & \multicolumn{2}{p{0.3\linewidth}}{
        \cellcolor[rgb]{0.9,0.9,0.9}{
            \makecell[{{p{\linewidth}}}]{
                \texttt{\tiny{[GM$|$GM]}}
                \texttt{HAMLE: C, (7, 7)} \\
            }
        }
    }
    & & \\ \\

    \theutterance \stepcounter{utterance}  
    & & & \multicolumn{2}{p{0.3\linewidth}}{
        \cellcolor[rgb]{0.9,0.9,0.9}{
            \makecell[{{p{\linewidth}}}]{
                \texttt{\tiny{[GM$|$GM]}}
                \texttt{C moved to (7, 7)} \\
            }
        }
    }
    & & \\ \\

    \theutterance \stepcounter{utterance}  
    & & \multicolumn{4}{p{0.6\linewidth}}{
        \cellcolor[rgb]{0.9,0.9,0.9}{
            \makecell[{{p{\linewidth}}}]{
                \texttt{\tiny{[GM$\rangle$P2]}}
                \texttt{Mesajın diğer oyuncuya iletildi.} \\
\\ 
\\ 
\texttt{Şu anda en fazla 20 turdan 9. turu oynuyorsunuz.} \\
\texttt{Toplamda 12 cezanın 12 kadarını biriktirdiniz.Diğer oyuncu kendi ızgarasında bir nesneyi hareket ettirdi.} \\
\\ 
\\ 
\texttt{Bir sonraki komutun nedir?} \\
            }
        }
    }
    & \\ \\

    \theutterance \stepcounter{utterance}  
    & & & \multicolumn{4}{p{0.6\linewidth}}{
        \cellcolor[rgb]{0.9,0.9,0.9}{
            \makecell[{{p{\linewidth}}}]{
                \texttt{\tiny{[GM$\langle$P2]}}
                \texttt{HAMLE: W, (4, 5)} \\
            }
        }
    }
     \\ \\

    \theutterance \stepcounter{utterance}  
    & & & \multicolumn{2}{p{0.3\linewidth}}{
        \cellcolor[rgb]{0.9,0.9,0.9}{
            \makecell[{{p{\linewidth}}}]{
                \texttt{\tiny{[GM$|$GM]}}
                \texttt{HAMLE: W, (4, 5)} \\
            }
        }
    }
    & & \\ \\

    \theutterance \stepcounter{utterance}  
    & & & \multicolumn{2}{p{0.3\linewidth}}{
        \cellcolor[rgb]{0.9,0.9,0.9}{
            \makecell[{{p{\linewidth}}}]{
                \texttt{\tiny{[GM$|$GM]}}
                \texttt{W moved to (4, 5)} \\
            }
        }
    }
    & & \\ \\

    \theutterance \stepcounter{utterance}  
    & & \multicolumn{4}{p{0.6\linewidth}}{
        \cellcolor[rgb]{0.9,0.9,0.9}{
            \makecell[{{p{\linewidth}}}]{
                \texttt{\tiny{[P1$\langle$GM]}}
                \texttt{C’ başarıyla (7,7) koordinatlarına taşındı. Güncellenmiş ızgaran şu şekilde:} \\
\texttt{```} \\
\texttt{ 1234567} \\
\texttt{╔════╤══╗} \\
\texttt{║◌◌◌T│◌◌║ 1} \\
\texttt{║◌◌◌◌│◌◌║ 2} \\
\texttt{║◌◌◌◌│◌◌║ 3} \\
\texttt{║H◌◌◌│◌◌║ 4} \\
\texttt{╟──┐W│◌I║ 5} \\
\texttt{║◌◌├─┼──╢ 6} \\
\texttt{║◌◌│◌│◌C║ 7} \\
\texttt{╚══╧═╧══╝} \\
\texttt{```} \\
\texttt{Şu anda en fazla 20 turdan 9. turu oynuyorsunuz.} \\
\texttt{Toplamda 12 cezanın 12 kadarını biriktirdiniz.Diğer oyuncu kendi ızgarasında bir nesneyi hareket ettirdi.} \\
\\ 
\\ 
\texttt{Bir sonraki komutun nedir?} \\
            }
        }
    }
    & \\ \\

    \theutterance \stepcounter{utterance}  
    & \multicolumn{4}{p{0.6\linewidth}}{
        \cellcolor[rgb]{0.9,0.9,0.9}{
            \makecell[{{p{\linewidth}}}]{
                \texttt{\tiny{[P1$\rangle$GM]}}
                \texttt{HAMLE: I, (2, 4)} \\
            }
        }
    }
    & & \\ \\

    \theutterance \stepcounter{utterance}  
    & & & \multicolumn{2}{p{0.3\linewidth}}{
        \cellcolor[rgb]{0.9,0.9,0.9}{
            \makecell[{{p{\linewidth}}}]{
                \texttt{\tiny{[GM$|$GM]}}
                \texttt{HAMLE: I, (2, 4)} \\
            }
        }
    }
    & & \\ \\

    \theutterance \stepcounter{utterance}  
    & & & \multicolumn{2}{p{0.3\linewidth}}{
        \cellcolor[rgb]{0.9,0.9,0.9}{
            \makecell[{{p{\linewidth}}}]{
                \texttt{\tiny{[GM$|$GM]}}
                \texttt{I moved to (2, 4)} \\
            }
        }
    }
    & & \\ \\

    \theutterance \stepcounter{utterance}  
    & & \multicolumn{4}{p{0.6\linewidth}}{
        \cellcolor[rgb]{0.9,0.9,0.9}{
            \makecell[{{p{\linewidth}}}]{
                \texttt{\tiny{[GM$\rangle$P2]}}
                \texttt{W’ başarıyla (4,5) koordinatlarına taşındı. Güncellenmiş ızgaran şu şekilde:} \\
\texttt{```} \\
\texttt{ 1234567} \\
\texttt{╔════╤══╗} \\
\texttt{║◌◌◌T│◌◌║ 1} \\
\texttt{║◌◌◌◌│◌◌║ 2} \\
\texttt{║◌◌◌◌│◌◌║ 3} \\
\texttt{║H◌◌◌│◌◌║ 4} \\
\texttt{╟──┐W│◌◌║ 5} \\
\texttt{║◌◌├─┼──╢ 6} \\
\texttt{║I◌│◌│◌C║ 7} \\
\texttt{╚══╧═╧══╝} \\
\texttt{```} \\
\texttt{Şu anda en fazla 20 turdan 10. turu oynuyorsunuz.} \\
\texttt{Toplamda 12 cezanın 12 kadarını biriktirdiniz.Diğer oyuncu kendi ızgarasında bir nesneyi hareket ettirdi.} \\
\\ 
\\ 
\texttt{Bir sonraki komutun nedir?} \\
            }
        }
    }
    & \\ \\

    \theutterance \stepcounter{utterance}  
    & & & \multicolumn{4}{p{0.6\linewidth}}{
        \cellcolor[rgb]{0.9,0.9,0.9}{
            \makecell[{{p{\linewidth}}}]{
                \texttt{\tiny{[GM$\langle$P2]}}
                \texttt{HAMLE: I, (2, 4)} \\
            }
        }
    }
     \\ \\

    \theutterance \stepcounter{utterance}  
    & & & \multicolumn{2}{p{0.3\linewidth}}{
        \cellcolor[rgb]{0.9,0.9,0.9}{
            \makecell[{{p{\linewidth}}}]{
                \texttt{\tiny{[GM$|$GM]}}
                \texttt{HAMLE: I, (2, 4)} \\
            }
        }
    }
    & & \\ \\

    \theutterance \stepcounter{utterance}  
    & & & \multicolumn{2}{p{0.3\linewidth}}{
        \cellcolor[rgb]{0.9,0.9,0.9}{
            \makecell[{{p{\linewidth}}}]{
                \texttt{\tiny{[GM$|$GM]}}
                \texttt{I moved to (2, 4)} \\
            }
        }
    }
    & & \\ \\

    \theutterance \stepcounter{utterance}  
    & & \multicolumn{4}{p{0.6\linewidth}}{
        \cellcolor[rgb]{0.9,0.9,0.9}{
            \makecell[{{p{\linewidth}}}]{
                \texttt{\tiny{[P1$\langle$GM]}}
                \texttt{I’ başarıyla (2,4) koordinatlarına taşındı. Güncellenmiş ızgaran şu şekilde:} \\
\texttt{```} \\
\texttt{ 1234567} \\
\texttt{╔════╤══╗} \\
\texttt{║◌◌◌T│◌◌║ 1} \\
\texttt{║◌◌◌◌│◌◌║ 2} \\
\texttt{║◌◌◌◌│◌◌║ 3} \\
\texttt{║HI◌◌│◌◌║ 4} \\
\texttt{╟──┐W│◌◌║ 5} \\
\texttt{║◌◌├─┼──╢ 6} \\
\texttt{║◌◌│◌│◌C║ 7} \\
\texttt{╚══╧═╧══╝} \\
\texttt{```} \\
\texttt{Şu anda en fazla 20 turdan 10. turu oynuyorsunuz.} \\
\texttt{Toplamda 12 cezanın 12 kadarını biriktirdiniz.Diğer oyuncu kendi ızgarasında bir nesneyi hareket ettirdi.} \\
\\ 
\\ 
\texttt{Bir sonraki komutun nedir?} \\
            }
        }
    }
    & \\ \\

    \theutterance \stepcounter{utterance}  
    & \multicolumn{4}{p{0.6\linewidth}}{
        \cellcolor[rgb]{0.9,0.9,0.9}{
            \makecell[{{p{\linewidth}}}]{
                \texttt{\tiny{[P1$\rangle$GM]}}
                \texttt{SÖYLE: Benim tahtamda C(7,7), I(2,4), H(1,4), T(4,1) hedeflerinde. W şu an (4,5) {-} hedefte. Tüm nesneler hedef koordinatlarda. Senin durumun nasıl? Hepsi yerindeyse "bitti mi?" diye sorayım.} \\
            }
        }
    }
    & & \\ \\

    \theutterance \stepcounter{utterance}  
    & & & \multicolumn{2}{p{0.3\linewidth}}{
        \cellcolor[rgb]{0.9,0.9,0.9}{
            \makecell[{{p{\linewidth}}}]{
                \texttt{\tiny{[GM$|$GM]}}
                \texttt{SÖYLE: Benim tahtamda C(7,7), I(2,4), H(1,4), T(4,1) hedeflerinde. W şu an (4,5) {-} hedefte. Tüm nesneler hedef koordinatlarda. Senin durumun nasıl? Hepsi yerindeyse "bitti mi?" diye sorayım.} \\
            }
        }
    }
    & & \\ \\

    \theutterance \stepcounter{utterance}  
    & & \multicolumn{4}{p{0.6\linewidth}}{
        \cellcolor[rgb]{0.9,0.9,0.9}{
            \makecell[{{p{\linewidth}}}]{
                \texttt{\tiny{[GM$\rangle$P2]}}
                \texttt{I’ başarıyla (2,4) koordinatlarına taşındı. Güncellenmiş ızgaran şu şekilde:} \\
\texttt{```} \\
\texttt{ 1234567} \\
\texttt{╔════╤══╗} \\
\texttt{║◌◌◌T│◌◌║ 1} \\
\texttt{║◌◌◌◌│◌◌║ 2} \\
\texttt{║◌◌◌◌│◌◌║ 3} \\
\texttt{║HI◌◌│◌◌║ 4} \\
\texttt{╟──┐W│◌◌║ 5} \\
\texttt{║◌◌├─┼──╢ 6} \\
\texttt{║◌◌│◌│◌C║ 7} \\
\texttt{╚══╧═╧══╝} \\
\texttt{```} \\
\texttt{Şu anda en fazla 20 turdan 11. turu oynuyorsunuz.} \\
\texttt{Toplamda 12 cezanın 12 kadarını biriktirdiniz.Diğer oyuncu şu mesajı gönderdi:} \\
\\ 
\\ 
\texttt{"Benim tahtamda C(7,7), I(2,4), H(1,4), T(4,1) hedeflerinde. W şu an (4,5) {-} hedefte. Tüm nesneler hedef koordinatlarda. Senin durumun nasıl? Hepsi yerindeyse "bitti mi?" diye sorayım."} \\
\\ 
\\ 
\texttt{Bir sonraki komutun nedir?} \\
            }
        }
    }
    & \\ \\

    \theutterance \stepcounter{utterance}  
    & & & \multicolumn{4}{p{0.6\linewidth}}{
        \cellcolor[rgb]{0.9,0.9,0.9}{
            \makecell[{{p{\linewidth}}}]{
                \texttt{\tiny{[GM$\langle$P2]}}
                \texttt{HAMLE: C, (7, 7)} \\
            }
        }
    }
     \\ \\

    \theutterance \stepcounter{utterance}  
    & & & \multicolumn{2}{p{0.3\linewidth}}{
        \cellcolor[rgb]{0.9,0.9,0.9}{
            \makecell[{{p{\linewidth}}}]{
                \texttt{\tiny{[GM$|$GM]}}
                \texttt{HAMLE: C, (7, 7)} \\
            }
        }
    }
    & & \\ \\

    \theutterance \stepcounter{utterance}  
    & & & \multicolumn{2}{p{0.3\linewidth}}{
        \cellcolor[rgb]{0.9,0.9,0.9}{
            \makecell[{{p{\linewidth}}}]{
                \texttt{\tiny{[GM$|$GM]}}
                \texttt{Ceza: (7,7) boş değil; "C" içeriyor.} \\
\\ 
\texttt{Toplamda 12 cezanın 13 kadarını biriktirdiniz.Lütfen tekrar deneyin!} \\
            }
        }
    }
    & & \\ \\

    \theutterance \stepcounter{utterance}  
    & & \multicolumn{4}{p{0.6\linewidth}}{
        \cellcolor[rgb]{0.9,0.9,0.9}{
            \makecell[{{p{\linewidth}}}]{
                \texttt{\tiny{[GM$\rangle$P2]}}
                \texttt{Ceza: (7,7) boş değil; "C" içeriyor.} \\
\\ 
\texttt{Toplamda 12 cezanın 13 kadarını biriktirdiniz.Lütfen tekrar deneyin!} \\
            }
        }
    }
    & \\ \\

    \theutterance \stepcounter{utterance}  
    & & & \multicolumn{4}{p{0.6\linewidth}}{
        \cellcolor[rgb]{0.9,0.9,0.9}{
            \makecell[{{p{\linewidth}}}]{
                \texttt{\tiny{[GM$\langle$P2]}}
                \texttt{SÖYLE: C zaten (7,7) konumundaydı. Tüm nesneler hedef koordinatlarda: H(1,4), T(4,1), W(4,5), I(2,4), C(7,7). Ben bitti. Sende bitti mi?} \\
            }
        }
    }
     \\ \\

    \theutterance \stepcounter{utterance}  
    & & & \multicolumn{2}{p{0.3\linewidth}}{
        \cellcolor[rgb]{0.9,0.9,0.9}{
            \makecell[{{p{\linewidth}}}]{
                \texttt{\tiny{[GM$|$GM]}}
                \texttt{SÖYLE: C zaten (7,7) konumundaydı. Tüm nesneler hedef koordinatlarda: H(1,4), T(4,1), W(4,5), I(2,4), C(7,7). Ben bitti. Sende bitti mi?} \\
            }
        }
    }
    & & \\ \\

    \theutterance \stepcounter{utterance}  
    & & \multicolumn{4}{p{0.6\linewidth}}{
        \cellcolor[rgb]{0.9,0.9,0.9}{
            \makecell[{{p{\linewidth}}}]{
                \texttt{\tiny{[P1$\langle$GM]}}
                \texttt{Mesajın diğer oyuncuya iletildi.} \\
\\ 
\\ 
\texttt{Şu anda en fazla 20 turdan 11. turu oynuyorsunuz.} \\
\texttt{Toplamda 12 cezanın 13 kadarını biriktirdiniz.Diğer oyuncu şu mesajı gönderdi:} \\
\\ 
\\ 
\texttt{"C zaten (7,7) konumundaydı. Tüm nesneler hedef koordinatlarda: H(1,4), T(4,1), W(4,5), I(2,4), C(7,7). Ben bitti. Sende bitti mi?"} \\
\\ 
\\ 
\texttt{Bir sonraki komutun nedir?} \\
            }
        }
    }
    & \\ \\

    \theutterance \stepcounter{utterance}  
    & \multicolumn{4}{p{0.6\linewidth}}{
        \cellcolor[rgb]{0.9,0.9,0.9}{
            \makecell[{{p{\linewidth}}}]{
                \texttt{\tiny{[P1$\rangle$GM]}}
                \texttt{SÖYLE: bitti!} \\
            }
        }
    }
    & & \\ \\

    \theutterance \stepcounter{utterance}  
    & & & \multicolumn{2}{p{0.3\linewidth}}{
        \cellcolor[rgb]{0.9,0.9,0.9}{
            \makecell[{{p{\linewidth}}}]{
                \texttt{\tiny{[GM$|$GM]}}
                \texttt{SÖYLE: bitti!} \\
            }
        }
    }
    & & \\ \\

    \theutterance \stepcounter{utterance}  
    & & & \multicolumn{2}{p{0.3\linewidth}}{
        \cellcolor[rgb]{0.9,0.9,0.9}{
            \makecell[{{p{\linewidth}}}]{
                \texttt{\tiny{[GM$|$GM]}}
                \texttt{true} \\
            }
        }
    }
    & & \\ \\

    \theutterance \stepcounter{utterance}  
    & & & \multicolumn{2}{p{0.3\linewidth}}{
        \cellcolor[rgb]{0.9,0.9,0.9}{
            \makecell[{{p{\linewidth}}}]{
                \texttt{\tiny{[GM$|$GM]}}
                \texttt{Initial states:} \\
\texttt{```} \\
\texttt{Player 1:} \\
\texttt{ 1234567} \\
\texttt{╔════╤══╗} \\
\texttt{║H◌◌◌│◌◌║ 1} \\
\texttt{║◌W◌◌│◌◌║ 2} \\
\texttt{║◌C◌◌│◌◌║ 3} \\
\texttt{║T◌◌◌│I◌║ 4} \\
\texttt{╟──┐◌│◌◌║ 5} \\
\texttt{║◌◌├─┼──╢ 6} \\
\texttt{║◌◌│◌│◌◌║ 7} \\
\texttt{╚══╧═╧══╝} \\
\\ 
\\ 
\texttt{Player 2:} \\
\texttt{ 1234567} \\
\texttt{╔════╤══╗} \\
\texttt{║◌◌◌T│◌◌║ 1} \\
\texttt{║◌◌◌◌│◌◌║ 2} \\
\texttt{║◌◌◌◌│◌◌║ 3} \\
\texttt{║W◌◌◌│◌◌║ 4} \\
\texttt{╟──┐◌│◌◌║ 5} \\
\texttt{║H◌├─┼──╢ 6} \\
\texttt{║I◌│◌│◌C║ 7} \\
\texttt{╚══╧═╧══╝} \\
\\ 
\texttt{```} \\
            }
        }
    }
    & & \\ \\

    \theutterance \stepcounter{utterance}  
    & & & \multicolumn{2}{p{0.3\linewidth}}{
        \cellcolor[rgb]{0.9,0.9,0.9}{
            \makecell[{{p{\linewidth}}}]{
                \texttt{\tiny{[GM$|$GM]}}
                \texttt{End states:} \\
\texttt{```} \\
\texttt{Player 1:} \\
\texttt{ 1234567} \\
\texttt{╔════╤══╗} \\
\texttt{║◌◌◌T│◌◌║ 1} \\
\texttt{║◌◌◌◌│◌◌║ 2} \\
\texttt{║◌◌◌◌│◌◌║ 3} \\
\texttt{║HI◌◌│◌◌║ 4} \\
\texttt{╟──┐W│◌◌║ 5} \\
\texttt{║◌◌├─┼──╢ 6} \\
\texttt{║◌◌│◌│◌C║ 7} \\
\texttt{╚══╧═╧══╝} \\
\\ 
\\ 
\texttt{Player 2:} \\
\texttt{ 1234567} \\
\texttt{╔════╤══╗} \\
\texttt{║◌◌◌T│◌◌║ 1} \\
\texttt{║◌◌◌◌│◌◌║ 2} \\
\texttt{║◌◌◌◌│◌◌║ 3} \\
\texttt{║HI◌◌│◌◌║ 4} \\
\texttt{╟──┐W│◌◌║ 5} \\
\texttt{║◌◌├─┼──╢ 6} \\
\texttt{║◌◌│◌│◌C║ 7} \\
\texttt{╚══╧═╧══╝} \\
\\ 
\texttt{```} \\
            }
        }
    }
    & & \\ \\

    \theutterance \stepcounter{utterance}  
    & & & \multicolumn{2}{p{0.3\linewidth}}{
        \cellcolor[rgb]{0.9,0.9,0.9}{
            \makecell[{{p{\linewidth}}}]{
                \texttt{\tiny{[GM$|$GM]}}
                \texttt{* success: True} \\
\texttt{* lose: False} \\
\texttt{* aborted: False} \\
\texttt{{-}{-}{-}{-}{-}{-}{-}} \\
\texttt{* Shifts: 6.00} \\
\texttt{* Max Shifts: 8.00} \\
\texttt{* Min Shifts: 4.00} \\
\texttt{* End Distance Sum: 0.00} \\
\texttt{* Init Distance Sum: 23.71} \\
\texttt{* Expected Distance Sum: 20.95} \\
\texttt{* Penalties: 13.00} \\
\texttt{* Max Penalties: 12.00} \\
\texttt{* Rounds: 11.00} \\
\texttt{* Max Rounds: 20.00} \\
\texttt{* Object Count: 5.00} \\
            }
        }
    }
    & & \\ \\

\end{supertabular}
}

\end{document}
